\scanpage{3}
$k$ corps de car $p>0$

$G$ grp $p$-divisible, $G^*$ dual
\[
g=\dim G,\quad g^*=\dim G^*
\]
$\widetilde G$ grp \unsure{enveloppant}
\diagram{\begin{tikzcd}
0\arrow[r]&\underline V(t_{G^*})\arrow[r]&\widetilde G\arrow[r]&G\arrow[r]&0
\end{tikzcd}}
\diagram{\begin{tikzcd}[column sep=small,row sep=large]
0\arrow[r]&\underset{\omega_G=}{t_G^\vee}\arrow[r]&H(G)\arrow[r]\arrow[d,shift right,"H(V)=V_0"']&t_{G^*}\arrow[r]&0\\
0\arrow[r]&\underset{\omega_G^{(p)}=}{t_G^{\vee(p)}}\arrow[r]&H(G)^{(p)}\arrow[r]\arrow[u,shift right,"H(F)=F_0"']&t_{G^*}^{(p)}\arrow[r]&0\\
&\text{dim }g&\text{dim }g+g^*&\text{dim }g^*
\end{tikzcd}\quad\vrule height 4cm depth 0cm\quad
\begin{tikzcd}[column sep=small,row sep=large]
\text{duales} & & & & \\
0\arrow[r]&\underset{\omega_{G^*}=}{t_{G^*}^\vee}\arrow[r]&t_{\widetilde G}\arrow[r]\arrow[d,shift right,"V_0"']&t_G\arrow[r]&0\\
0\arrow[r]&\underset{\omega_{G^*}^{(p)}=}{t_{G^*}^{\vee(p)}}\arrow[r]&t_{\widetilde G}^{(p)}\arrow[r]\arrow[u,shift right,"F_0"']&t_G^{(p)}\arrow[r]&0
\end{tikzcd}}
\[
[\ \uline{\text{NB}}\quad H(G)=t_{\unsure{\widetilde{G^*}|G^*}}\deleted{\unsure{\vee}}\ ]
\]
\textcircled{1} \unsure{suites sont} pour $t_{\widetilde G}$ et $t_{\widetilde{G^*}}$ transposées l'une de l'autre pour un accouplement canonique fonctoriel, \deleted{\unsure{avec}} pour lequel $F_0$ et $V_0$ (si on est en car $p$) s'interchangent.

\textcircled{2} Dans la suite exacte de gauche, on a
\[
\begin{aligned}
\operatorname{Ker}F_0&=\operatorname{Im}V_0=t_G^{\vee(p)}=\omega_G^{(p)}\\
\operatorname{Ker}V_0&=\operatorname{Im}F_0=\text{\unsure{ce que ça peut} (\unsure{innovation intéressante})}
\end{aligned}
\]
\textcircled{3} $F_0$ et $V_0$ opérant sur $H(G)$ définissent le module de Dieudonné \deleted{de \unsure{$G$}} réduit mod $p$.

\margin{?}
\textcircled{4} La composée
\diagram{\begin{tikzcd}
t_{G^*}^{(p)}\arrow[r,"F_0"]&H(G)\arrow[r,dash]&t_{G^*}
\end{tikzcd}}
\deleted{\unsure{la multiplication}\ \illegible{2}}
\margin{correspond : la puissance $p$-ième dans \unsure{l'algèbre} [et correspond : Frobenius dans $H^{\unsure{2|1}}(G,\mathbf G_a)$]\quad $H(G)=t_{\widetilde{G^*}}$}

\uline{Corollaire}
\[
\begin{array}{rcl}
G\text{ purement infinitésimal}&\Longleftrightarrow&F_0\text{ nilpotent}\qquad\text{i.e. }G\text{ sans tore}\\
G^*\text{ purement infinitésimal}&\Longleftrightarrow&V_0\text{ nilpotent}\\[-2pt]
\text{i.e. }G\text{ sans composante}\ \text{toroïdale}&&\\
G\text{ et }G^*\text{ purement infinitésimaux}&\Longleftrightarrow&F_0\text{ et }V_0\text{ nilpotents}
\end{array}
\]
\mspage{2}
